\section{雇用と職場：専門能力を発揮できる条件を守る}
\label{theory:employment:background}
技術者は安全な製品を作る責任だけでなく、設計・製造・点検を担う人の尊厳と健康にも関わる。威圧によって異常を報告できない職場では、技術上の問題を発見しても対処へ結び付かない。長時間労働や特定個人への業務集中も判断・継承を損なう。雇用法令と職場の仕組みを、専門能力が継続して発揮されるための条件として理解する。

\subsection{労働法：主体と義務を先に読む}
\label{theory:employment:law}
伝統的な労働三法は労働基準法・労働組合法・労働関係調整法である\parencite{employment_labor_three}。最低条件の保護、集団的な交渉、紛争の調整という役割を区別する。労働安全衛生法は労働者の安全と健康を扱い、規模等に応じた産業医の選任を求める\parencite{employment_industrial_physician}。技術士資格があるだけで、労働者を使う際の義務が免除されるわけではない。

36協定は法定の時間外・休日労働に関する労使協定である。締結と届出を区別し、届出だけで上限を自由に撤廃できるものと読まない\parencite{employment_overtime}。以下では出題当時の記述を保ち、2026年10月12日時点で確認した変更は別に述べる。全ての告示の改正履歴を照合し終えたという意味ではない。

\subsection{パワハラ：職位より関係、叱責より態様}
\label{theory:employment:power}
平成24年WG報告は、職場内の優位性を背景に、適正な業務範囲を超えて苦痛又は環境悪化を生む行為を整理した。上司から部下だけでなく、知識や協力への依存による同僚・部下の優位性も含む。身体攻撃・精神攻撃・隔離・過大要求・過小要求・私生活への過度な介入は典型例であり、網羅的な六項目の暗記だけでは足りない\parencite{employment_power2012}。

検討では、目的に必要な指示か、方法が相当か、相手と周囲の環境へどう影響するかを分ける。手順の誤りを具体的に指摘することと、人前で人格を否定することは同じでない。主観的不満だけで該当性は決まらないが、相手の苦痛を調べずに「指導だから問題なし」としてもいけない。現在の指針も三要素と客観的な適正指導の区別を定める\parencite{employment_power_current}。

\subsection{セクハラ：不利益と環境悪化}
\label{theory:employment:sexual}
性的言動への対応で処遇上の不利益を受ける対価型と、言動によって仕事の環境が害される環境型を区別する。同性間でも生じ、労働者の雇用形態や行為者の所属だけで保護を除外できない。派遣労働者には派遣元・派遣先の関係を確認する\parencite{employment_sexual_current}。

拒絶を口にできない関係や、問題化すると仕事を失うと感じる状況があるため、沈黙だけから同意を推定しない。相談を受けた組織は事実を確認し、情報の扱いと相談者への不利益を防ぐ体制を整える。2018年研修資料も措置義務を説明しており、単なる努力目標という過去問の肢は誤る\parencite{employment_sexual2018}。

\subsection{各類型の範囲を同一にしない}
\label{theory:employment:scope}
日常語のハラスメントは広い。セクハラ、パワハラ、妊娠等に関するハラスメントは法的な条件が同じではない。顧客からの性的言動もセクハラの行為者範囲に入るが、その事実だけで顧客との全ての紛争を同じ職場内のパワハラ定義へ入れない。古いWG報告の顧客例は、職場内の定義とは別の脚注である。

性的指向や性自認を侮辱したり、本人の同意なく周囲へ伝えたりすることは、人格と機微情報に関わる。技術指導に必要な行為かを確認する。過去問の「本人はもとより職場全体のハラスメントとなり得る」は、周囲にも就業環境の悪化が及び得るという判断であり、全員が必ず同一の法的被害者となるという意味ではない。

\subsection{産休と妊婦への配慮}
\label{theory:employment:maternity}
産前産後休業は母体保護を目的とし、育児休業とは根拠と条件が異なる。産前は原則出産予定前6週間、多胎14週間で本人の請求が条件である。産後は原則8週間の就業禁止で、6週間経過後に本人が請求し医師が支障なしと認める業務には例外がある。正規・有期・派遣といった雇用形態で一律に除外しない\parencite{employment_maternity2014,employment_maternity_current}。

妊娠を理由に仕事から一方的に排除することと、健康や業務状況を考えて本人へ調整を相談することは異なる。配慮という名称だけで不利益を正当化せず、目的・本人の意向・法定措置を確認する。2018年の公的研修例も、残業の負担を相談して減らす配慮と、権利利用を妨げる強要を分けている。

\subsection{労働時間：場所と自己管理は免除理由でない}
\label{theory:employment:worktime}
使用者の指揮命令下の時間を考える。命令された業務上の学習は、学習という名称でも労働時間に含まれる。使用者は時間を適切に把握し、自己申告を用いる場合にも記録との食違いを確認する\parencite{employment_hours2017}。

管理監督者は肩書だけで決まらず、職務・権限・勤務の裁量・待遇等の実態から判断する。該当する場合も深夜割増賃金や年次有給休暇まで一律に除外されるものではない\parencite{employment_manager2008}。勤務終了後に必要な報告や研修を行わせながら記録から除外することは、健康管理と制度の趣旨を損なう。

\subsection{裁量労働：資格名と業務の列挙を分ける}
\label{theory:employment:discretion}
専門業務型裁量労働制は、法令で指定する業務の中から労使協定で対象を定め、定めた時間働いたものとみなす仕組みである。平成25年当時は19業務であり、保存した平成26年3月資料にも同じ19業務が並ぶ。この一覧に「技術士の業務」はない\parencite{employment_discretion2014}。

技術士が研究開発など対象業務に従事する可能性と、技術士資格を持てば自動的に対象になる命題は異なる。この混同で平成25年度問4の初回解答を誤ったため、初回3と資料確認後の訂正5を別々に保存した。適用を確かめるには肩書から推測せず、実際の従事業務と指定範囲を照合する。

\paragraph{現在の変更}2024年4月から対象は20業務で、本人同意・同意の撤回等の要件が加わった。現行一覧にも技術士業務という列挙はない\parencite{employment_discretion_current}。みなし時間は健康管理のための勤務状況把握を不要にするものではない。

\subsection{フレックスとテレワーク}
\label{theory:employment:flex}
フレックスは一定の清算期間の総労働時間を定め、その枠内で始業・終業を労働者が選択する仕組みである。期間の上限は2019年4月の改正で1か月から3か月へ延長された\parencite{employment_reform2018}。平成29年度の説明へ3か月を遡及しない。柔軟な時間配分と、総労働時間の把握を両立する。

\label{theory:employment:telework}
テレワークは場所の選択を広げるが、雇用型では労働法令が適用される。勤務場所を示し、時間を把握し、中抜けで自由利用が保障される時間と、移動中も命令された仕事をする時間を分ける。平成30年指針の概要を載せた令和2年白書でもこの区別を確認できる\parencite{employment_telework2020}。時間単位年休には別の制度要件があり、中抜けしたことだけで勝手に年休へ振り替えない。

現在の指針は記録・申告の方法、休止と移動、長時間化対策を整理する\parencite{employment_telework_current}。通信できる環境は常時応答の義務と同じでない。自己管理という説明によって、組織側の時間把握や負荷調整を放棄しない。

\subsection{働き方改革：上限・割増・休暇を分ける}
\label{theory:employment:reform}
法定の時間外労働上限は使用者の指示だけで緩和できない。年10日以上の年休が付与される労働者について、年5日の取得を確保する。本人の請求による取得や計画付与の取得分を含むため、全員へ別枠の5日を追加指定する説明は不正確である\parencite{employment_paidleave}。

月60時間超の時間外労働は割増賃金の問題であり、10時間ごとに代休1日という制度ではない。代替休暇には労使協定と本人意思を要し、引き上げた割増分に対応して換算する仕組みである\parencite{employment_alternative2008}。法定休日に働いた後の代休、事前に休日を振り替える操作、割増分の代替休暇は別の概念である。

\subsection{仕事と生活の調和：属人化と不要作業を減らす}
\label{theory:employment:worklife}
業務を標準化して引き継げる形にし、資料の過剰作成や探索を減らし、予定と負荷を共有する。これらは同じ仕事を少人数へ集中させるためでなく、時間と相互支援を確保する組織改善である。内閣府の十の実践は会議、資料、整理、標準化、時間、分担、相互理解、予定、集中、改善共有を扱う\parencite{employment_worklife2010}。

会議の到達点を明らかにする場合、技術上の証拠が不十分なら追加調査と責任者・期限を決めることも到達点になる。「必ず結論」という標語を、判断の根拠を省略して安全性を断定する命令へ変えない。効率化の取組後も、業務の質、負荷の偏り、実際の労働時間を観察し、改善の効果を確認する。
